% Entwurf Abschnitt 6.3 (AP5), Stand 20.11.2026. Code-/Doku-Stand: G1@25d5c1a; Pfade in "% Basis:" relativ zu g1pilot/.
% Darstellungen (fig:betriebszustaende, tab:haltstufen, tab:schutzmechanismen, fig:armuebergabe) baut AP1/AP2;
% hier nur \ref und Erläuterung, keine Umgebungen. Grenzwerte stehen nur in tab:schutzmechanismen.
% Umfang Ziel: ca. 3 Seiten (6.3.1 ~1 S., 6.3.2 ~1,3 S., 6.3.3 ~0,7 S.).

\section{Sicherheitskonzept}
\label{sec:sicherheitskonzept}
% [A5, F3]
Das Sicherheitskonzept setzt die Anforderung A5 aus Tabelle~\ref{tab:anforderungen} um: Ein Bedienfehler oder der Ausfall eines Bedienwegs soll am realen Roboter nicht zu einer unbegrenzten Bewegung führen, und ein Halt soll den stehenden Roboter nicht zu Fall bringen.
Dazu legt Abschnitt~\ref{sec:zustaende} Betriebszustände und Stufen des Halts fest, Abschnitt~\ref{sec:schutzmechanismen} ordnet die Schutzmechanismen, und Abschnitt~\ref{sec:armuebergabe} behandelt den Wechsel der Armregelung zwischen Onboard-Regler und Umgebung.

\subsection{Betriebszustände und Übergänge}
\label{sec:zustaende}

% Art G: Leitgedanke "aktiver sicherer Zustand", Rückverweis 3.4.3
Ausgangspunkt des Konzepts ist, dass der sichere Zustand eines balancierenden Zweibeiners ein aktiv geregelter Zustand ist (Abschnitt~\ref{sec:sicherheit_zweibeiner}).
In der klassischen Maschinensicherheit entzieht ein Schutzstopp die Energie, und die Maschine kommt zu einem ungefährlichen Stillstand \cite{Ding2026}\ev{Ding2026.02}.
Beim dynamisch balancierenden Humanoiden führt derselbe Energieentzug dagegen zu einem unkontrollierten Sturz, der selbst eine Gefährdung ist \cite{Ding2026}\ev{Ding2026.03}; sein sicherer Zustand ist nach einem nicht begutachteten Preprint ein aktiv geregelter Zustand \cite{Ding2026}\ev{Ding2026.01}.
Auch mit gebremsten und abgeschalteten Gelenken ist ein Humanoid nicht notwendig sicher \cite{Prather2025}\ev{Prather2025.03}.
Im Vergleich zu Vier- oder Sechsbeinern liegt sein Schwerpunkt relativ hoch und die konvexe Hülle seiner Füße ist klein, was ihn anfälliger für Stürze macht \cite{Subburaman2023}\ev{Subburaman2023.01}, und ein Sturz kann Roboter, Objekte und Personen schädigen \cite{Subburaman2023}\ev{Subburaman2023.07}.

% Art G (Ding2026.08/.09, DINENISO13850.02) + Art E (Folgerung für die Umgebung)
% Basis: G1@25d5c1a docs/51-T-Notaus.md, Abschnitt „Wer was tut“; docs/50-A-Echter-Roboter.md, Abschnitt „Notaus-Reflex“
Daraus folgt ein Halt in Stufen statt eines einzigen Abschaltens.
Ein zu abrupter Halt kann eine Störung jenseits der Erholungsgrenze des Balancereglers einbringen und einen Sturz auslösen \cite{Ding2026}\ev{Ding2026.09}; die Norm zum Not-Halt verlangt dagegen, gefahrbringende Bewegungen ohne zusätzliche Gefährdungen anzuhalten \cite{DINENISO13850}\ev{DINENISO13850.02}.
In einer Fallstudie wird deshalb ein kontrollierter Halt im Stand befohlen, bei dem die Regelung zum Balancieren erhalten bleibt \cite{Ding2026}\ev{Ding2026.08}.
Die Umgebung trennt nach demselben Gedanken zwei Wirkungen: Der Not-Halt entlastet die Arme, von denen bei Kontakt die unmittelbare Gefährdung ausgeht, und lässt die Beine unter der Regelung des Onboard-Reglers stehen.

% Art E: abstrahiertes Zustandsdiagramm (N7: Nils „vermutlich ja“)
% Basis: G1@25d5c1a docs/30-T-Stehen-und-Gehen.md, Abschnitte „Zustandsmaschine“ und „loco_client (echter Roboter)“;
%        docs/20-T-Arm-Manipulation.md, Abschnitt „Zustandsmaschine“; docs/51-A-Notaus.md, Abschnitte „Stufen“, „In der Simulation“, „Quittieren“;
%        Code g1pilot/navigation/loco_sim.py (HOLD/STAND/WALK/DAMP, _on_start -> HOLD, Sturzerkennung fall_gz/fall_debounce_s),
%        g1pilot/navigation/loco_client.py (FSM 1/4/500/501, _waist_blocks_walk, Bodenkontakt knee_load_min_nm),
%        g1pilot/teleoperation/status_line.py (start_buttons: Rückfrage bei Damp und Walk), gui_base.py (einziger Sender von /g1pilot/start)
Abbildung~\ref{fig:betriebszustaende} fasst die Betriebszustände in einem Diagramm zusammen, das für Simulation und Realbetrieb gilt.
Es abstrahiert vier gekoppelte Zustandsautomaten, nämlich Onboard-Regler bzw.\ Ersatzregler, Armregler und Bedienmodus, zu den Zuständen \emph{Gedämpft}, \emph{Bereit} (am Kran hängend bzw.\ in der Simulation am Rumpf gehalten) und \emph{Stehen} mit genau einem der Bedienmodi GREIFEN und GEHEN.
\qf{Herstellerdokumentation G1: Zustände des Onboard-Reglers (gedämpft, vorbereitet, Gehen), Rückverweis sec:g1\_onboard; Benutzerhandbuch G1 nicht in Zotero (L16)}
Das Gehen setzt im Realbetrieb voraus, dass der Roboter balanciert, die Hüftdrehung in der Mitte steht und die Knie Last tragen; Dämpfen und Gehen verlangen eine zweite Bestätigung.
Der Not-Halt ist aus jedem Zustand erreichbar und verriegelt alle Bewegungsbefehle, bis an der Bedienoberfläche oder am Smartphone quittiert wird, bewusst nicht am Gamepad.
Danach behalten die Beine im Realbetrieb ihren Zustand, die Arme bleiben schlaff, bis die Bedienperson sie im Modus GREIFEN bewusst wieder übernimmt.

% Art E: Halt-Stufen
% Basis: G1@25d5c1a docs/51-A-Notaus.md, Abschnitte „Stufen“, „In der Simulation“; docs/51-T-Notaus.md, Abschnitt „Wer was tut“;
%        Code g1pilot/manipulation/arm_controller.py (_on_emergency_stop, _on_full_emergency_stop, _publish_arm_slack),
%        g1pilot/navigation/loco_client.py (_halt_legs: Geschwindigkeit 0 statt Damp; Damp nur mit estop_damp_legs, Default false)
Tabelle~\ref{tab:haltstufen} stellt die fünf Stufen des Halts mit ihrer Wirkung auf Arme, Hände, Taille und Beine gegenüber; sie greifen von oben nach unten tiefer ein.
\emph{Stopp} beendet die laufende Armbewegung, \emph{Gehen stoppen} hält den Roboter beim Loslassen des Totmanns im Stand an.
Im zentralen \emph{Not-Halt} sinken die Arme drehmomentfrei und gedämpft ab, Hände und Taille halten ihre Stellung, und die Beine erhalten nur einen Haltebefehl, sodass der Onboard-Regler den Roboter stehend weiter balanciert.
Erst der \emph{Voll-Notaus} dämpft auch die Beine, sodass der Roboter zusammensackt und am Kran gesichert sein muss; die \emph{Fernbedienung des Herstellers} wirkt als letzte Stufe an der Firmware.
In der Simulation dämpft bereits der Not-Halt die Beine, und eine erkannte Sturzneigung führt dort automatisch in den Zustand Gedämpft.

% Art G (DGUV2015.10) + Art E: Analogie mit genanntem Unterschied
% Kann
Die Stufen entsprechen den Stoppkategorien aus Abschnitt~\ref{sec:schutzfunktionen} nur als Analogie.
Kategorie~0 setzt durch sofortiges Abschalten der Energie still, Kategorie~2 gesteuert und mit beibehaltener Energie \cite{DGUV2015}\ev{DGUV2015.10}.
Der Not-Halt ähnelt für die Beine der Kategorie~2, endet aber nicht in einem Stillstand, sondern in einem aktiv gehaltenen Gleichgewicht; der Voll-Notaus kommt der Kategorie~0 am nächsten, dämpft die Gelenke jedoch, statt die Energie abzuschalten.
Ob Zustände und Übergänge wie beschrieben greifen, weist Abschnitt~\ref{sec:eval_sicherheit} nach.

\subsection{Schutzmechanismen}
\label{sec:schutzmechanismen}

% Art E: betrachtete Gefährdungen (eigene Zusammenstellung, keine Gefährdungsbeurteilung, N9) und Wirkprinzipien
% Basis: G1@25d5c1a docs/50-A-Echter-Roboter.md, Abschnitt „E · Guarantees, die bereits im Code stehen“; docs/51-A-Notaus.md
Der Not-Halt allein genügt für A5 nicht, weil er erst wirkt, nachdem eine Person die Gefährdung erkannt hat.
Betrachtet werden Stoß und Quetschen durch einen Arm, Kollision mit der Umgebung oder dem eigenen Körper, Sturz, unbeabsichtigtes Loslaufen sowie der Ausfall einer Verbindung oder der Bedienperson.
Die Umgebung staffelt ihre Mechanismen deshalb nach vier Wirkprinzipien: \emph{Begrenzen}, \emph{Vermeiden}, \emph{Überwachen} und \emph{Anhalten}.
\pruefen{Begründung für die Staffelung nach vier Wirkprinzipien von Nils bestätigen (in der G1-Doku nicht als Begründung dokumentiert)}
Tabelle~\ref{tab:schutzmechanismen} ordnet ihnen die Mechanismen mit Gefährdung, Auslegungswert, Geltung in Simulation und Realbetrieb und dem Abschnitt der Umsetzung zu.

% Art E + Art G (Orientierung Handgeschwindigkeit)
% Basis: G1@25d5c1a docs/50-A-Echter-Roboter.md, Abschnitt E („Harte Grenzen nur auf real“, „Sanftes Anfahren und begrenzter Schleppfehler“);
%        docs/20-T-Arm-Manipulation.md, Abschnitte „main_loop() – pro Tick“ und „Übergabe an den Onboard-Regler“;
%        Code g1pilot/manipulation/arm_controller.py (REAL_MAX_GAINS, REAL_MAX_ARM_VEL, REAL_MAX_EE_VEL, _apply_real_safety_limits;
%        Kommentare Z. 167/1586/3888 „reduced speed nach ISO 10218-1 / ISO TS 15066“), launch/bringup_real.launch.py Z. 97–111
\emph{Begrenzen} wirkt auf jede Armbewegung, gleich von welchem Bedienweg das Ziel kommt.
Begrenzt sind Gelenkgeschwindigkeit und -beschleunigung der Arme, die Geschwindigkeit der Hand, der Abstand zwischen Sollwert und Messwert, die Reglerverstärkungen und die Gehgeschwindigkeit; der begrenzte Abstand verhindert, dass sich an einem festgehaltenen Arm ein Sprung aufbaut.
Die harten Grenzen gelten nur im Realbetrieb, abweichende Vorgaben der Konfiguration korrigiert die Umgebung beim Start.
Die Grenze der Handgeschwindigkeit orientiert sich an der reduzierten Geschwindigkeit, mit der Industrieroboter in der Betriebsart T1 verfahren \cite{DGUV2015}\ev{DGUV2015.04}, und an der Grundannahme einer sicher reduzierten Geschwindigkeit in ISO/TS~15066 \cite{Villani2018}\ev{Villani2018.09}; anders als dort ist sie keine sicher überwachte Geschwindigkeit, sondern eine Sollwertbegrenzung in Software.

% Art E: Vermeiden
% Basis: G1@25d5c1a docs/20-T-Arm-Manipulation.md, Abschnitt „Laufzeit-Gate“; docs/30-T-Stehen-und-Gehen.md, Abschnitt „loco_client“ (Bodenkontakt vor »Walk«);
%        docs/50-A-Echter-Roboter.md, Abschnitt E („Kein Auto-Start“, „WALK ist gesperrt, bis balanciert“); start.sh Z. 189–207 (Bestätigung „REAL“)
\emph{Vermeiden} setzt vor der Bewegung an.
Eine Kollisionsprüfung zur Laufzeit lässt keine weitere Annäherung der Arme an die Umgebung unter einen Mindestabstand zu und hält bei drohender Selbstkollision an; gespeicherte Posen werden über geplante Bahnen angefahren (Abschnitt~\ref{sec:kollision_sdt}).
Hinzu kommen die Verriegelungen des Gehens aus Abschnitt~\ref{sec:zustaende}, die Bestätigung beim Start des Realbetriebs und der Grundsatz, dass im Realbetrieb ohne Eingabe keine Bewegung beginnt.

% Art E + Art G (DGUV2015.06/.07): Totmann als Analogie zur Zustimmeinrichtung
% Basis: G1@25d5c1a docs/51-A-Notaus.md, Abschnitt „Was zusätzlich von selbst anhält“; docs/51-T-Notaus.md, Abschnitt „Zeitgrenzen ohne Notaus-Taste“;
%        docs/11-T-Controller.md, Abschnitt „Sicherheit“ (Lebendprüfung der Funkverbindung); docs/12-T-Handy.md, Abschnitt „Zeiten“;
%        Code g1pilot/teleoperation/ps4_walk.py (Docstring: L2 gibt frei, Sticks geben vor; DEADMAN_THRESHOLD 0.5 = gehalten, keine Panikstellung)
\emph{Überwachen} hält an, sobald eine Voraussetzung der Bewegung entfällt.
Das Gehen per Gamepad ist nur freigegeben, solange eine Totmanntaste gehalten wird; die Bewegungsvorgabe kommt von den Steuerknüppeln.
Die Taste ähnelt damit einer Zustimmeinrichtung, die nur bei ständiger Betätigung eine Maschinenfunktion erlaubt und allein keine Bewegung starten darf \cite{DGUV2015}\ev{DGUV2015.06}.
Anders als die für Industrieroboter vorgeschriebenen dreistufigen Zustimmungsschalter \cite{DGUV2015}\ev{DGUV2015.07} gibt sie das Gehen aber auch bei vollständigem Durchdrücken frei und wird nicht sicherheitsgerichtet ausgewertet.
Zeitüberwachungen beenden das Gehen, wenn Gamepad-Signal oder Gehbefehle ausbleiben, eine Lebendprüfung erkennt eine eingefrorene Funkverbindung des Gamepads, und das Smartphone sperrt ohne aktuellen Status alles außer dem Not-Halt.

% Art E + Art G (DGUV2015.09, DINENISO13850.01): Anhalten
% Basis: G1@25d5c1a docs/51-A-Notaus.md, Abschnitte „Stufen“, „Quittieren“, „Was zusätzlich von selbst anhält“ (Handy-NOT-HALT läuft über den PC);
%        docs/51-T-Notaus.md, Abschnitt „Parameter“ (G1_PHONE_VIEW_ESTOP Default 1)
\emph{Anhalten} umfasst die Halt-Stufen aus Abschnitt~\ref{sec:zustaende}.
Der Not-Halt ist an Bedienoberfläche, Gamepad und Smartphone auslösbar, am Smartphone auch in der Rolle „Ansicht“ und am Gamepad auch im gesperrten Zustand.
Damit orientiert sich die Umgebung an der Forderung, dass die Not-Halt-Funktion jederzeit verfügbar ist und Vorrang vor allen anderen Funktionen hat \cite{DINENISO13850}\ev{DINENISO13850.01}.
Anders als die für Industrieroboter vorgeschriebene Not-Halt-Einrichtung \cite{DGUV2015}\ev{DGUV2015.09} ist er aber eine Nachricht in der Software; der Not-Halt des Smartphones läuft über den PC und wirkt nicht, wenn dieser hängt.

% Art G (Ding2026.05, Prather2025.06) + Art E: Einordnung, unabhängige Ebene (N9: keine Aussage zu einer Gefährdungsbeurteilung)
% Basis: G1@25d5c1a docs/50-A-Echter-Roboter.md, Abschnitte „A · Vor dem Einschalten“ (zweite Person mit Fernbedienung, Kran/Traggurt) und „Notaus-Reflex“
Alle Mechanismen sind Softwarefunktionen auf einem handelsüblichen PC ohne sicherheitstechnische Bewertung.
Auch die Bordrechner des G1 sind handelsübliche Standardhardware, sodass ein darauf laufender Sicherheitscontroller keine zertifizierte Sicherheitslaufzeit ist \cite{Ding2026}\ev{Ding2026.05}, und die bestehenden Normen lassen sich nicht direkt auf Humanoide anwenden \cite{Prather2025}\ev{Prather2025.06}.
Die Mechanismen begrenzen deshalb, was ein Bedienfehler auslösen kann, ersetzen aber keine unabhängige Schutzebene.
Diese bilden im Realbetrieb eine zweite Person mit der Fernbedienung des Herstellers, die unabhängig von PC und Umgebung wirkt, und die Sicherung am Kran oder Traggurt bei den ersten Geh- und Balanceversuchen.
\qf{Herstellerangaben zur Fernbedienung (Dämpfung als Not-Halt, Wirkung parallel zu eigener Software); Benutzerhandbuch G1 nicht in Zotero (L23)}
Den Nachweis der Mechanismen führt Abschnitt~\ref{sec:eval_sicherheit}.

\subsection{Übergabe der Arme an den Onboard-Regler}
\label{sec:armuebergabe}

% Art E (eigene Beobachtung) + Rückverweis 2.3.3; Herstellerbeleg fehlt (L24)
% Basis: G1@25d5c1a docs/20-T-Arm-Manipulation.md, Abschnitte „rt/arm_sdk-Gewichtsrampe“ und „Übergabe an den Onboard-Regler“
Im Realbetrieb regelt der Onboard-Regler auch die Arme, solange die Umgebung sie nicht übernimmt; die Herstellerschnittstelle überblendet dazu über ein Gewicht zwischen beiden Sollwerten (Abschnitt~\ref{sec:g1_armsdk}).
\qf{Herstellerdokumentation zur Überblendung per Gewicht (arm\_sdk) und ihrer Wirkung auf die Taille; nicht in Zotero (L24)}
Jeder Wechsel des Gewichts zieht den Arm von seiner Stellung in die Pose des anderen Reglers, sodass eine reine Gewichtsrampe einen weit vorn stehenden Arm nach unten schwingen lässt.
Statt einer einheitlichen Rampe legt die Umgebung deshalb für vier Fälle je eine Regel fest; Abbildung~\ref{fig:armuebergabe} zeigt den Verlauf von Gewicht und Armstellung für die ersten beiden schematisch.

% Art E: vier Fälle
% Basis: G1@25d5c1a docs/50-A-Echter-Roboter.md, Abschnitt E („Übernehmen ohne Bewegung“, „Zurückgeben ohne Schwung“, „Programmende“);
%        docs/20-T-Arm-Manipulation.md, Abschnitte „Übergabe an den Onboard-Regler“, „E-Stop-Semantik“;
%        Code g1pilot/manipulation/arm_controller.py (_record_onboard_pose, _start_release, _finish_release, shutdown_release, _on_emergency_stop)
Beim \emph{Übernehmen} kommandiert der Armregler genau die gemessene Armstellung des Onboard-Reglers und erhöht erst dann das Gewicht über eine Rampe, sodass sich der Arm nicht bewegt.
Beim \emph{Zurückgeben} fährt der Arm zuerst geplant und langsam in diese gemerkte Pose, erst dort sinkt das Gewicht; währenddessen werden neue Ziele abgelehnt, und nach einer Zeitgrenze folgt die Rampe aus der erreichten Stellung.
Im \emph{Not-Halt} bleibt das Gewicht bei eins und die Gelenke erhalten nur Dämpfung, denn bei Gewicht null führe der Onboard-Regler die Arme aktiv in seine Grundpose.
Am \emph{Programmende} fahren die Arme außerhalb des Not-Halts kurz in Richtung der gemerkten Pose, bevor die Rampe einsetzt, statt dass die Befehle abreißen.
% Kann
In der Simulation schaltet das Gewicht ohne Rampe um, die Regeln betreffen daher vor allem den Realbetrieb.

% Art E + N6 (Realtest als Begründung, ohne Datum)
% Basis: G1@25d5c1a docs/50-A-Echter-Roboter.md, Abschnitt E („Taille wird mitgehalten“); docs/51-T-Notaus.md, Abschnitt „Wer was tut“;
%        Code g1pilot/manipulation/arm_controller.py (_publish_arm_slack schreibt Taille, estop_waist_hold Default true), launch/bringup_real.launch.py Z. 128.
%        Hinweis: docs/20-T-Arm-Manipulation.md „Bekannte Einschränkungen“ sagt noch, die Taille werde nie beschrieben – veraltet, Code gilt.
Das Gewicht nimmt beim G1 nach eigener Beobachtung auch die Taille aus der Regelung des Onboard-Reglers.
\pruefen{Wirkung des Gewichts auf die Taille als eigene Beobachtung bestätigen, solange kein Herstellerbeleg vorliegt (L24)}
Die Umgebung hält die Taille deshalb selbst in ihrer gemessenen Stellung, auch im Not-Halt; den Ausschlag gab ein Realtest, in dem eine schlaffe Taille den Oberkörper des balancierenden Roboters kippen ließ und der Roboter loslief \pruefen{Doku des Realtests von Nils}.
Die Übergaben weisen Abschnitt~\ref{sec:eval_realtests} am Roboter und Abschnitt~\ref{sec:eval_sicherheit} in automatisierten Tests nach.
